\subsection{Thermoelastic Wave Propagation in Geological Media}

Thermoelastic wave propagation in geological media can be understood as the propagation of a mechanical disturbance through a deformable solid whose stress state is influenced by temperature change. In a purely elastic medium, waves arise because particles are displaced from equilibrium, producing strain and stress. Spatial variations in stress then cause further motion, allowing the disturbance to travel through the material. In a thermoelastic medium, this mechanical process is modified by the fact that temperature perturbations produce thermal strain and thermal stress. Therefore, the wave field is controlled not only by density and elastic stiffness, but also by the thermal state of the rock and by the coupling between temperature and deformation \cite{aki2002,nowacki1986,boley1960}.

Using the coupled thermoelastic framework introduced in the previous section, the mechanical part of wave propagation is described by the equation of motion
\begin{equation}
\label{eq:eom}
\rho \frac{\partial^2 u_i}{\partial t^2}
=
\frac{\partial \sigma_{ij}}{\partial x_j}.
\end{equation}
This equation states that stress gradients produce acceleration of material particles. In the context of wave propagation, a local disturbance creates a changing displacement field; this displacement field produces strain; strain produces stress; and spatial variation in stress drives the motion forward through the medium. Thus, displacement, strain, stress, and acceleration evolve together during elastic wave propagation \cite{aki2002,jaeger2007}.

The thermoelastic nature of the process enters through the stress relation
\begin{equation}
\sigma_{ij}
=
\lambda \delta_{ij} \varepsilon_{kk}
+
2\mu \varepsilon_{ij}
-
\gamma \delta_{ij} \theta .
\end{equation}
The first two terms describe the ordinary elastic stress produced by mechanical deformation, while the final term represents the stress contribution associated with temperature perturbation. This means that heating or cooling can modify the stress field through which the wave propagates. If temperature changes are spatially non-uniform or mechanically constrained, they may produce additional stress gradients and therefore influence the motion of the medium \cite{nowacki1986,boley1960}.

The thermal field is coupled to the mechanical field through the heat-conduction equation with a thermoelastic coupling term:
\begin{equation}
k \nabla^2 T
-
C \frac{\partial T}{\partial t}
=
\gamma T_0 \frac{\partial \varepsilon_{kk}}{\partial t}.
\end{equation}
This equation shows that temperature evolution is governed by heat conduction and by changes in volumetric deformation. The term involving \(\partial \varepsilon_{kk}/\partial t\) is important because it links the thermal field to compressional deformation. Therefore, thermoelastic wave propagation is not simply the superposition of heat conduction and elastic waves; it is a coupled process in which temperature affects stress and deformation affects temperature \cite{nowacki1986,boley1960}.

In this context, P-waves and S-waves have different thermoelastic significance. P-waves are compressional waves and involve changes in volume. Since the thermoelastic coupling term depends on volumetric strain \(\varepsilon_{kk}\), compressional motion is directly connected with temperature perturbations in the classical coupled theory. S-waves mainly involve shear deformation and do not produce volume change in the same way in an ideal isotropic medium. For this reason, direct thermoelastic coupling is usually more closely associated with compressional deformation than with purely shear deformation. In real rocks, however, fractures, pores, anisotropy, and heterogeneity may complicate this distinction because shear motion can interact with cracks, pore fluids, and non-uniform stress fields \cite{aki2002,mavko2009,jaeger2007}.

Temperature-induced stress perturbations may also modify the mechanical state of the medium. If heating or cooling is non-uniform, different parts of the rock expand or contract by different amounts. This produces local thermal stress and may change the effective stiffness, crack state, and wave velocity of the material. Such effects are especially relevant near geothermal systems, cooling igneous bodies, underground openings, or fractured zones where heat transfer is uneven. In these settings, the propagation of mechanical disturbances cannot be interpreted independently of the temperature field \cite{robertson1988,turcotte2014,jaeger2007}.

Another important feature of thermoelastic wave propagation is attenuation. Heat conduction is a diffusive process, whereas elastic wave propagation is mechanical and dynamic. When deformation produces local temperature changes, heat may flow from warmer regions to cooler regions. This irreversible redistribution of thermal energy can convert part of the mechanical energy of the wave into heat, reducing the amplitude of the disturbance. In natural rocks, attenuation is also affected by cracks, pore fluids, friction, scattering, and mineral heterogeneity, but thermoelastic heat conduction provides one classical mechanism for dissipation \cite{boley1960,nowacki1986,mavko2009}.

Material boundaries influence thermoelastic wave propagation because they introduce contrasts in physical properties. When an elastic wave reaches an interface between rocks with different density or elastic stiffness, part of the energy may be reflected and part may be transmitted or refracted. In a thermoelastic setting, contrasts in thermal expansion coefficient, thermal conductivity, and heat capacity can additionally influence the local thermal stress field near the boundary. Thus, lithological contacts are not only elastic contrasts but may also act as thermoelastic contrasts \cite{aki2002,mavko2009,robertson1988}.

Geological heterogeneity makes thermoelastic wave propagation spatially variable. Natural rocks contain minerals with different elastic and thermal properties, pores with different shapes, fractures with different orientations, and layers with different stiffness and conductivity. Therefore, the values of \(\rho\), \(\lambda\), \(\mu\), \(k\), \(C\), and \(\gamma\) may vary within the medium. Elastic disturbances may accelerate or decelerate, scatter or change direction as they encounter material contrasts, while thermal disturbances may travel faster through one part of a rock than through another. The thermoelastic response is therefore controlled not only by the governing equations, but also by the geological structure of the medium \cite{mavko2009,jaeger2007,robertson1988}.

Layered, fractured, porous, and anisotropic rocks may show direction-dependent propagation. In sedimentary rocks, bedding may cause waves traveling parallel and perpendicular to layers to behave differently. In fractured rock masses, aligned cracks or joints may reduce stiffness more strongly in some directions than in others. In porous rocks, pore shape and saturation may influence both wave velocity and heat transfer. In crystalline rocks, mineral orientation and microcrack patterns may also introduce directional behavior. Since thermoelastic propagation depends on both elastic and thermal properties, direction dependence may arise from anisotropy in stiffness, thermal conductivity, crack orientation, or geological fabric \cite{mavko2009,jaeger2007,aki2002}.

Thermoelastic wave propagation in geological media should therefore be understood as a coupled thermal--mechanical process occurring within a structured natural material. Elastic waves are governed by inertia and stress gradients, while temperature perturbations modify the stress state through thermal expansion. Heat conduction controls the evolution of temperature and can contribute to wave attenuation. In ideal homogeneous isotropic media, this behavior can be described by classical thermoelastic equations. In real rocks, mineral composition, porosity, fractures, bedding, saturation, and anisotropy make the response spatially variable and direction-dependent \cite{nowacki1986,boley1960,mavko2009,turcotte2014}.